
\documentclass{article}
\usepackage{colortbl}
\usepackage{makecell}
\usepackage{multirow}
\usepackage{supertabular}

\begin{document}

\newcounter{utterance}

\centering \large Interaction Transcript for game `codenames', experiment `risk\_high', episode 4 with mistral{-}large{-}3{-}azure.
\vspace{24pt}

{ \footnotesize  \setcounter{utterance}{1}
\setlength{\tabcolsep}{0pt}
\begin{supertabular}{c@{$\;$}|p{.15\linewidth}@{}p{.15\linewidth}p{.15\linewidth}p{.15\linewidth}p{.15\linewidth}p{.15\linewidth}}
    \# & $\;$A & \multicolumn{4}{c}{Game Master} & $\;\:$B\\
    \hline

    \theutterance \stepcounter{utterance}  
    & & & \multicolumn{2}{p{0.3\linewidth}}{
        \cellcolor[rgb]{0.9,0.9,0.9}{
            \makecell[{{p{\linewidth}}}]{
                \texttt{\tiny{[GM$|$GM]}}
                {'hidden': {'team': ['valgus', 'haigus', 'pall', 'oja', 'kell', 'vibu', 'leping', 'post', 'mexico'], 'innocent': ['kaubik', 'nupp', 'riik'], 'opponent': ['haigla', 'apelsin', 'mudel', 'sild', 'õhk', 'tilk', 'tool', 'rist'], 'assassin': ['teemant', 'märge', 'šokolaad', 'berliin', 'kalju']}, 'revealed': {'team': {'team': [], 'innocent': [], 'opponent': [], 'assassin': []}, 'opponent': {'team': [], 'innocent': [], 'opponent': [], 'assassin': []}}}
            }
        }
    }
    & & \\ \\

    \theutterance \stepcounter{utterance}  
    & & \multicolumn{4}{p{0.6\linewidth}}{
        \cellcolor[rgb]{0.9,0.9,0.9}{
            \makecell[{{p{\linewidth}}}]{
                \texttt{\tiny{[P1$\langle$GM]}}
                \texttt{Esitage üks ühe sõnaga vihje, mis on seotud mõne teie meeskonna sõnaga allpool esitatud nimekirjast. Võite valida, kui paljude ja milliste sõnadega teie vihje seotud on.} \\
\texttt{Eesmärk on, et teie partner peab teie antud vihjete põhjal ära arvama kõik teie meeskonna sõnad, enne kui teine meeskond on kõik oma sõnad avastanud. Kui teie partner arvab ära "vastase" või "segaja" sõna, lõpeb nende arvamisvoor ja vastasmeeskond saab teha käigu. Kui teie partner arvab ära "tapa" sõna, lõpeb mäng ja teie meeskond kaotab kohe.} \\
\texttt{Mida rohkemate sõnadega teie vihje seostub, seda kiiremini võite mängu võita, kuid seda raskem võib olla teie partneril seoseid ära arvata, seega valige oma vihjesõna targalt. Vihjesõna peab olema semantiliselt seotud sihtsõnadega, see ei tohi olla üks nimekirjades olevatest sõnadest ega sisaldada nende osasid.} \\
\\ 
\texttt{Andke alati oma ühe sõnaga vihje ja oma komadega eraldatud sihtsõnade loetelu järgmises vormis ning tehke oma vastused võimalikult lühikeseks, ärge kunagi lisage muud teksti kui selles vormis nõutud:} \\
\\ 
\texttt{VIHJE: <SÕNA>} \\
\texttt{SIHID: <SIHID>} \\
\\ 
\texttt{Teie meeskonna sõnad on: valgus, haigus, pall, oja, kell, vibu, leping, post, mexico.} \\
\texttt{Teie vastase meeskonna sõnad on: haigla, apelsin, mudel, sild, õhk, tilk, tool, rist.} \\
\texttt{Segajate sõnad on: kaubik, nupp, riik.} \\
\texttt{Tapa sõnad on: teemant, märge, šokolaad, berliin, kalju.} \\
            }
        }
    }
    & \\ \\

    \theutterance \stepcounter{utterance}  
    & \multicolumn{4}{p{0.6\linewidth}}{
        \cellcolor[rgb]{0.9,0.9,0.9}{
            \makecell[{{p{\linewidth}}}]{
                \texttt{\tiny{[P1$\rangle$GM]}}
                \texttt{VIHJE: aeg} \\
\texttt{SIHID: kell, leping} \\
\\ 
\texttt{VIHJE: sport} \\
\texttt{SIHID: pall, vibu} \\
\\ 
\texttt{VIHJE: vool} \\
\texttt{SIHID: valgus, oja} \\
\\ 
\texttt{VIHJE: kiri} \\
\texttt{SIHID: post, mexico} \\
\\ 
\texttt{VIHJE: nakkus} \\
\texttt{SIHID: haigus} \\
            }
        }
    }
    & & \\ \\

    \theutterance \stepcounter{utterance}  
    & & & \multicolumn{2}{p{0.3\linewidth}}{
        \cellcolor[rgb]{0.9,0.9,0.9}{
            \makecell[{{p{\linewidth}}}]{
                \texttt{\tiny{[GM$|$GM]}}
                {'player': 'Cluegiver', 'type': 'rambling error', 'utterance': 'VIHJE: aeg\nSIHID: kell, leping\n\nVIHJE: sport\nSIHID: pall, vibu\n\nVIHJE: vool\nSIHID: valgus, oja\n\nVIHJE: kiri\nSIHID: post, mexico\n\nVIHJE: nakkus\nSIHID: haigus', 'message': 'Your answer contained more than two lines, please only give one clue and your targets on two separate lines.'}
            }
        }
    }
    & & \\ \\

\end{supertabular}
}

\end{document}
